Let the first result be $r$. The five results are $r,\dfrac{3}{2}r,\dfrac{9}{4}r,\dfrac{27}{8}r,\dfrac{81}{16}r$. All are integers, so $16\mid r$. Since all are products from the $1$ through $10$ multiplication table, the last is at most $100$, forcing $r=16$. The fourth result is $54$, whose only factor pair with both factors at most $10$ is $6,9$. Thus Stepa multiplied $\boxed6\text{ and }\boxed9$.
